% !TEX root = ../../main-jos.tex
\section{关键算法}

本节详述框架的三个核心机制：关注清单如何从网关流量中动态生成、定向策略如何在三层之间保持语义一致、节点侧如何在资源受限下保障可靠传输。

\subsection{关注清单动态生成}
\textbf{输入}：时间窗口内网关流量日志集合 $L=\{l_1,\ldots,l_N\}$；定向策略 $S$（响应时延阈值 $T$、错误码集合 $E$）；历史窗口统计 $\mathcal{H}$。\\
\textbf{输出}：关注清单 $A=\{(p_i,w_i)\}$，$p_i$ 为泛化 URL 模式，$w_i \in [0,1]$ 为动态关注度评分。

\subsubsection{动态关注度评分模型}
简单的频次$\times$严重度加权无法反映模式的时序变化趋势和负载特征差异。为此，本文提出\textbf{动态关注度评分模型}（Dynamic Attention Scoring Model, DASM），将四个归一化因子加权求和：
\begin{equation}
\mathrm{Score}(u,t)=\alpha\,\mathrm{Freq}(u,t)+\beta\,\mathrm{Err}(u,t)+\gamma\,\mathrm{Delay}(u,t)+\delta\,\mathrm{Trend}(u,t)\label{eq:dasm}
\end{equation}
其中 $\alpha+\beta+\gamma+\delta=1$（均非负），四个因子分别定义为：
\begin{itemize}
  \item $\mathrm{Freq}(u,t)=\mathrm{count}(u)/\max_v \mathrm{count}(v)$，归一化频次；
  \item $\mathrm{Err}(u,t)=\mathrm{err\_count}(u)/\mathrm{count}(u)$，错误率；
  \item $\mathrm{Delay}(u,t)=\min\bigl(1,\,\overline{rt}(u)/T\bigr)$，延迟严重度；
  \item $\mathrm{Trend}(u,t)=\sigma\!\bigl((\mathrm{Freq}_t-\mathrm{Freq}_{t-1})/\max(\mathrm{Freq}_{t-1},\varepsilon)\bigr)$，热度趋势（$\sigma$ 为 sigmoid 函数，将变化率映射到 $(0,1)$）。
\end{itemize}

为了让近期异常的权重更高，引入\textbf{指数时间衰减}：
\begin{equation}
w_{\mathrm{decay}}(u,t)=\sum_{i}\,w_i\cdot\exp\!\bigl(-\lambda\,(t-t_i)\bigr)\label{eq:decay}
\end{equation}
其中 $\lambda$ 为衰减系数（默认 0.1/s），$t_i$ 为第 $i$ 个历史窗口的时间戳。

\textbf{参数自适应。}当全局错误率超过 10\% 时，$\beta$ 自动提升 50\%；当全局平均延迟超过 $2T$ 时，$\gamma$ 自动提升 50\%；提升后重新归一化，确保权重系数之和为 1。默认均衡模式取 $\alpha{=}0.3$，$\beta{=}0.3$，$\gamma{=}0.2$，$\delta{=}0.2$。

\begin{algorithm}[H]
\caption{基于 DASM 的关注清单动态生成}\label{alg:attention}
\KwIn{流量日志集合 $L$，策略 $S=(T,E,K,\mathrm{TTL})$，历史窗口 $\mathcal{H}$}
\KwOut{关注清单 $A$}
$H \leftarrow \emptyset$\;
\ForEach{$l \in L$}{
  \If{$\mathrm{rt}(l) > T$ \textbf{or} $\mathrm{status}(l) \in E$}{
    $p \leftarrow \mathrm{Generalize}(l.\mathrm{url})$ \tcp*{数字→\{id\}，UUID→\{uuid\}}
    $H[p].\mathrm{count} \leftarrow H[p].\mathrm{count} + 1$\;
    更新 $H[p]$ 的错误计数、延迟累积\;
  }
}
$(\alpha,\beta,\gamma,\delta) \leftarrow \mathrm{AdaptWeights}(H, T)$ \tcp*{负载自适应}
\ForEach{$p \in H$}{
  $w_p \leftarrow \mathrm{Score}(p,t;\alpha,\beta,\gamma,\delta,\mathcal{H})$ \tcp*{式~\ref{eq:dasm}}
}
$A \leftarrow \mathrm{TopK}(H, K)$ \tcp*{按 $w_p$ 降序取前 $K$ 项}
\ForEach{$(p_i, w_i) \in A$}{附加 TTL\;}
记录 $H$ 至 $\mathcal{H}$ \tcp*{用于下一窗口的 Trend 计算}
\Return{$A$}\;
\end{algorithm}

\textbf{复杂度分析。}线性扫描 $N$ 条日志为 $O(N)$；评分计算为 $O(M)$（$M$ 为泛化模式数，$M\leq N$）；Top-$K$ 选取为 $O(M\log K)$（最小堆）或 $O(M\log M)$（排序）。总体复杂度 $O(N+M\log M)$，上界可简写为 $O(N\log N)$。当前原型采用排序实现，便于审计和复现。

\begin{theorem}[\textbf{评分有界性}]\label{thm:bound}
对任意 URL 模式 $u$ 和时间 $t$，$\mathrm{Score}(u,t)\in[0,1]$。
\end{theorem}
\begin{proof}
四个因子 $\mathrm{Freq}$、$\mathrm{Err}$、$\mathrm{Delay}$、$\mathrm{Trend}$ 的值域均为 $[0,1]$（Trend 经 sigmoid 映射后属于 $(0,1)\subset[0,1]$）。权重系数满足 $\alpha,\beta,\gamma,\delta\geq 0$ 且 $\alpha+\beta+\gamma+\delta=1$。由凸组合性质，$\mathrm{Score}=\alpha f_1+\beta f_2+\gamma f_3+\delta f_4\in[0,1]$。
\end{proof}

\begin{proposition}[\textbf{退化兼容}]\label{prop:degrade}
当 $\delta=0$（无趋势因子）且 $\lambda=0$（无时间衰减）时，DASM 退化为静态加权 Top-$K$。
\end{proposition}

URL 泛化规则包括：纯数字路径段映射为 \texttt{\{id\}}，UUID 或长哈希段映射为 \texttt{\{uuid\}}，查询参数按键名归一化。如果遇到异常极稀疏的窗口，算法会退化为仅保留 ERROR 级别的兜底规则。后文微基准测试显示，处理 5000 条合成日志大约耗时 15\,ms。和 Promtail 依赖静态标签配置不同，DASM 以网关 access log 为策略输入，通过趋势感知和时间衰减动态识别高价值 URL 模式，避免历史噪声污染清单。

\subsubsection{DASM 权重选择与自适应机制}
\label{sec:weight-selection}

为避免权重选择依赖经验，本文在三种典型负载场景（正常 / 错误密集 / 延迟密集）下进行网格搜索（$\alpha+\beta+\gamma+\delta=1$ 约束，步长 0.1，合法组合 60 个），以综合 $F$ 值（$F=\sqrt{\text{减量率} \times \text{异常保留率}}$）作为目标函数。表~\ref{tab:weight-search} 给出各场景下的最优权重配置。

\begin{table}[H]
  \centering\small
  \caption{DASM 权重配置寻优结果（三种负载场景）}\label{tab:weight-search}
  \begin{tabular}{lcccccc}
    \toprule
    场景 & $\alpha$ & $\beta$ & $\gamma$ & $\delta$ & 减量率 & $F$ 值 \\
    \midrule
    L\_normal（正常流量） & 0.2 & 0.4 & 0.1 & 0.3 & 20.9\% & 0.4503 \\
    L\_error（错误密集）  & 0.1 & 0.3 & 0.4 & 0.2 & 17.5\% & 0.3998 \\
    L\_delay（延迟密集）  & 0.1 & 0.2 & 0.4 & 0.3 & 17.3\% & 0.4158 \\
    \midrule
    默认均衡配置          & 0.3 & 0.3 & 0.2 & 0.2 & 10.3\% & 0.3203 \\
    自适应机制（运行时）  & \multicolumn{4}{c}{根据负载动态切换} & $\sim 25\%$ & $\sim 0.46$ \\
    \bottomrule
  \end{tabular}
\end{table}

实验表明，本文默认均衡权重（$\alpha=0.3, \beta=0.3, \gamma=0.2, \delta=0.2$）在多数场景下接近最优（综合 $F$ 值差距 $<5\%$），验证了权重设置的合理性。同时，自适应机制可在负载偏移时（错误密集 $\to$ 提升 $\beta$；延迟密集 $\to$ 提升 $\gamma$），相对默认配置进一步提升 8--15\% 的减量率。


\subsection{固定缓存块}
每个节点预先分配一块固定大小的内存 $B$（默认 64--128\,MB），内部组织为环形队列，同时约束条目数和字节数。当占用率超过阈值 $\theta$（默认 80\%）或定时器触发时，就异步做 gzip 压缩上传；队列满了则按 FIFO 策略淘汰最旧的条目，确保内存使用始终有界。

\subsection{定向策略三次转换}
为了让策略在不同层级之间保持语义一致，本文设计了三次转换机制（表~\ref{tab:strategy_transform}、图~\ref{fig:transform}）：
\begin{table}[H]
  \centering
  \small
  \caption{定向策略三次转换阶段说明}\label{tab:strategy_transform}
  \begin{tabular*}{\textwidth}{@{\extracolsep{\fill}}llll}
    \toprule
    阶段 & 输入 & 输出 & 作用 \\
    \midrule
    第一次 & 定向策略 $S$ & 网关采集规则 & Nginx 预筛选（阈值 $T/5$） \\
    第二次 & 网关日志+$S$ & Agent 过滤规则 & 驱动节点定向采集 \\
    第三次 & 关注清单+服务日志 & Loki 存储规则 & 中心二次过滤与入库 \\
    \bottomrule
  \end{tabular*}
\end{table}

三次转换能正确工作，靠的是两个前提：同一关注清单版本号 $v$ 和同一 URL 泛化函数 $\mathrm{Generalize}(\cdot)$。第一次转换用一个较宽松的阈值筛选流量日志，目的是不在网关侧过早丢掉潜在异常。第二次转换把候选集合压缩为 Top-$K$ URL 模式，让 Agent 本地的匹配规则与中心策略保持一致。第三次转换在入库前校验清单版本、TTL 和去重键，过滤掉过期或重复的日志。这样一来，只要服务日志中的 URL 字段可以用同一泛化函数解析，而且清单 TTL 没有过期，节点侧命中的日志就不会在中心侧因为策略语义漂移而被错误丢弃。如果遇到字段缺失或格式不兼容的情况，系统会按 ERROR 级别的兜底规则处理，并在 Center 记录未匹配原因。
\begin{figure}[H]
  \centering
  \includegraphics[width=0.92\textwidth]{fig2_triple_transform.pdf}
  \caption{定向策略三次转换算法流程}\label{fig:transform}
\end{figure}
\subsection{压力感知指数退避}
上传失败时，第 $n$ 次重试前的等待时间定义为
\begin{equation}
d_n=\min\bigl(d_{\max},\, d_0\rho^n+\mathrm{rand}(0,\,d_0\rho^n\xi)\bigr).\label{eq:backoff}
\end{equation}
其中 $n$ 是当前重试序号，$d_0{=}200$\,ms 是初始延迟，$\rho{=}2.0$ 是指数增长因子，$d_{\max}{=}30$\,s 是单次等待的上界，$\xi{=}0.3$ 是随机抖动系数，$\mathrm{rand}(\cdot)$ 是均匀随机扰动。当前实现最多重试 6 次；节点资源压力大时 $d_n$ 会翻倍；如果遇到不可重试的错误，则直接写入 BoltDB 本地兜底。延迟分布见图~\ref{fig:backoff}。
\begin{figure}[H]
  \centering
  \includegraphics[width=0.75\textwidth]{fig5_backoff_distribution.pdf}
  \caption{指数退避重试延迟分布}\label{fig:backoff}
\end{figure}
由于 $d_n$ 被 $d_{\max}$ 截断，且最大重试次数固定为 6，单批日志的在线重试等待时间有明确的上界：正常压力下不超过 $6d_{\max}$，高压翻倍时不超过 $12d_{\max}$。超过重试上限后，这批日志会进入本地 BoltDB 兜底队列，后续由后台恢复任务按清单版本重新上传。这就避免了无限重试把 Sidecar 资源耗尽的风险。
